\section*{第 \ref{ch:b3:modules} 章——主理想整环上的
模}

\begin{solution}{exo:b3:modules:1}
(a) 若 $\Q = \Z q_1 + \dots + \Z q_k$，设 $d$ 为
各 $q_i$ 的公分母：每个线性组合都属于 $\frac1d\Z$，
但 $\frac1{2d} \notin \frac1d\Z$。矛盾。

(b) 无挠：由 $nq = 0$ 且 $n \neq 0$ 可得 $q = 0$
在 $\Q$ 中。不自由：任意两个非零有理数 $\frac ab, \frac cd$
满足非平凡关系 $(bc)\frac ab - (ad)\frac cd =
0$，故基至多含一个元素；若 $\Q \cong \Z$，
则 $\Q = \Z q$ 为循环群，但 $\frac q2 \notin \Z q$。（且 $\Q \neq
0$。）

(c) \cref{thm:b3:modules:structure} 假设\emph{\omterm{def:b3:modules:free}{有限
生成}}，而 (a) 否定了它：并无矛盾——相反 $\Q$
表明该假设在命题
“无挠 $\Rightarrow$ 自由”中是必要的。
\end{solution}

\begin{solution}{exo:b3:modules:2}
$360 = 2^3\cdot3^2\cdot5$。划分：$3$ 的划分为 $(3), (2,1),
(1,1,1)$；$2$ 的划分为 $(2), (1,1)$；$1$ 的划分为 $(1)$。因此共有 $3 \times 2
\times 1 = 6$ 个群。\omterm{thm:b3:modules:structure}{初等因子} $\to$
\omterm{thm:b3:modules:smith}{不变因子}：
\begin{center}
\begin{tabular}{l l}
$\Z/8 \times \Z/9 \times \Z/5$ & $\Z/360$\\
$\Z/8 \times \Z/3 \times \Z/3 \times \Z/5$ & $\Z/3 \times
\Z/120$\\
$\Z/4 \times \Z/2 \times \Z/9 \times \Z/5$ & $\Z/2 \times
\Z/180$\\
$\Z/4 \times \Z/2 \times \Z/3 \times \Z/3 \times \Z/5$ & $\Z/6
\times \Z/60$\\
$(\Z/2)^3 \times \Z/9 \times \Z/5$ & $\Z/2 \times \Z/2 \times
\Z/90$\\
$(\Z/2)^3 \times \Z/3 \times \Z/3 \times \Z/5$ & $\Z/2 \times
\Z/6 \times \Z/30$
\end{tabular}
\end{center}
（由\omterm{thm:b3:modules:structure}{初等因子}求\omterm{thm:b3:modules:smith}{不变因子}：最大的 $d_s$ 收集每个素数的
最高幂，依次向下。）
阶为 $p^5$ 的交换群个数等于 $5$ 的划分数，即 $7$。
\end{solution}

\begin{solution}{exo:b3:modules:3}
$B$：$D_1 = \gcd(2,4,6,8) = 2$；$D_2 = \abs{\det B} = \abs{16 -
24} = 8$。\omterm{thm:b3:modules:smith}{不变因子} $d_1 = 2$，$d_2 = 8/2 = 4$：史密斯
标准形为 $\operatorname{diag}(2, 4)$，且 $\Z^2/B\Z^2 \cong \Z/2\Z
\times \Z/4\Z$。

$C$：对角，但\emph{不是}史密斯标准形（$2 \nmid 3$）。$D_1 =
\gcd(2,3,12) = 1$；$D_2 = \gcd(2\cdot3,\, 2\cdot12,\, 3\cdot12) =
\gcd(6, 24, 36) = 6$；$D_3 = 72$。故 $d = (1, 6, 12)$，且
$\Z^3/C\Z^3 \cong \Z/6\Z\times\Z/12\Z$——与
中国剩余定理相符：$\Z/2\times\Z/3\times\Z/12 \cong \Z/6\times\Z/12$。
\end{solution}

\begin{solution}{exo:b3:modules:4}
写 $B = Q\,\operatorname{diag}(d_1, \dots, d_n)\,P$，其中 $Q, P
\in GL_n(\Z)$（\cref{thm:b3:modules:smith}；各 $d_i$ 皆非零，
由 $\det B \neq 0$ 可知）。于是 $\Z^n/B\Z^n \cong \Z^n/
\operatorname{diag}(d)\Z^n = \prod_i \Z/d_i\Z$（复合同构
$x \mapsto Q^{-1}x$ 将 $\Z^n$ 中的 $B\Z^n$ 映为
$\operatorname{diag}(d)P\Z^n = \operatorname{diag}(d)\Z^n$）。
其阶为 $\prod \abs{d_i} = \abs{\det
\operatorname{diag}(d)} = \abs{\det B}$，因 $\det Q, \det P = \pm
1$。对 $L = \Z(2,0) + \Z(1,3)$：$B = \bigl(\begin{smallmatrix}2
& 1\\ 0 & 3\end{smallmatrix}\bigr)$，$D_1 = 1$，$D_2 = 6$：$\Z^2/L
\cong \Z/6\Z$，其阶为 $\abs{\det B} = 6$。
\end{solution}

\begin{solution}{exo:b3:modules:5}
(a) 子模：已在\cref{def:b3:modules:torsion}中验证。若 $a(x +
T(M)) = 0$ 于 $M/T(M)$ 中，且 $a \neq 0$，则 $ax \in T(M)$：$bax
= 0$，存在 $b \neq 0$ 使之，且 $ba \neq 0$（整环），故 $x \in
T(M)$，于是该陪集为零。$M/T(M)$ 无挠。

(b) $(X, Y) \subseteq K[X,Y]$ 无挠（它是
整环 $K[X,Y]$ 作用于自身上的子模）。假设它是自由的，则任意两个
元素 $P, Q$ 满足 $Q\cdot P - P \cdot Q = 0$，
当 $P \ne Q$ 均非零时，这是非平凡关系，故基至多含一个元素：
$(X, Y) = (P)$ 为主\omterm{def:b3:rings:ideal}{理想}——与
\cref{exo:b3:rings:6}(a) 矛盾。该模无挠且\omterm{def:b3:modules:free}{有限生成}
（由 $X, Y$ 生成），但并非自由：在非\omterm{def:b3:rings:pidufd}{主理想整环} $K[X,Y]$ 上，
结构定理不成立。
\end{solution}

\begin{solution}{exo:b3:modules:6}
(a) 若 $\Z = 2\Z \oplus C$，则投影 $\Z \to \Z/2\Z$
限制为同构 $C \cong \Z/2\Z$：$C$ 应是
$\Z$ 的子群，其非零元 $x$ 满足 $2x \in C
\cap 2\Z = 0$。但 $\Z$ 无挠：$C = 0$，从而 $\Z =
2\Z$——矛盾。

(b) $A^n/M$ \omterm{def:b3:modules:free}{有限生成}且无挠，故自由
（\cref{thm:b3:modules:structure}）：$A^n/M \cong A^r$，基为
$f_1, \dots, f_r$。取 $y_i \in A^n$ 为 $f_i$ 的原像，
并令 $F = Ay_1 + \dots + Ay_r$。对每个 $x \in A^n$ 有 $\pi(x)
= \sum a_if_i$，故 $x - \sum a_iy_i \in M$：$A^n = M + F$。若
$\sum a_iy_i \in M$，对其作用 $\pi$ 得 $\sum a_if_i = 0$，
从而所有 $a_i = 0$（由基），故 $M \cap F = 0$。因此 $A^n = M \oplus
F$。
\end{solution}

\begin{solution}{exo:b3:modules:7}
\cref{exo:b3:modules:3} 中的约化是有效的：取
\[
L = \begin{pmatrix} 1 & 0\\ -3 & 1\end{pmatrix},
\qquad
R = \begin{pmatrix} 1 & 2\\ 0 & -1 \end{pmatrix},
\qquad
LBR = \begin{pmatrix} 2 & 0\\ 0 & 4\end{pmatrix}
\]
（行变换 $R_2 \leftarrow R_2 - 3R_1$，列变换
$C_2 \leftarrow C_2 - 2C_1$，再 $C_2 \leftarrow -C_2$）。令
$y = R^{-1}x$（是 $(\Z/20\Z)^2$ 上的双射，$R$
在 $\Z$ 上可逆），则同余方程组 $Bx \equiv b \pmod{20}$
等价于
\[
2y_1 \equiv b_1, \qquad 4y_2 \equiv -3b_1 + b_2 \pmod{20}.
\]
同余式 $ky \equiv c \pmod{20}$ 可解，当且仅当
$\gcd(k, 20) \mid c$：解存在，当且仅当 $2 \mid b_1$ 且 $4
\mid b_2 - 3b_1$，即\emph{$b_1$ 为偶数且 $b_2 \equiv 3b_1
\pmod 4$}。可解时共有 $2 \times 4 = 8$ 个
模 $20$ 的解。
\end{solution}

\begin{solution}{exo:b3:modules:8}
(a) 幂零的 $u$ 满足 $\mu_u = X^k$；\omterm{thm:b3:modules:structure}{初等因子}
为 $X^{k_1} \geq \dots$，每个若尔当块 $J_{k_i}(0)$ 对应
$4$ 的一个划分中的每个部分：
\begin{center}
\begin{tabular}{l l l}
划分 & \omterm{thm:b3:modules:jordan}{若尔当标准形} & \omterm{thm:b3:modules:smith}{不变因子} \\
$(4)$ & $J_4$ & $X^4$\\
$(3,1)$ & $J_3 \oplus J_1$ & $X,\ X^3$\\
$(2,2)$ & $J_2 \oplus J_2$ & $X^2,\ X^2$\\
$(2,1,1)$ & $J_2 \oplus J_1 \oplus J_1$ & $X,\ X,\ X^2$\\
$(1,1,1,1)$ & $0$ & $X, X, X, X$
\end{tabular}
\end{center}
(b) 取 $u = J_2\oplus J_2$，$v = J_2 \oplus J_1 \oplus J_1$：
两者都有 $\chi = X^4$，$\mu = X^2$，但\omterm{thm:b3:modules:smith}{不变因子}不同，
故不相似（见\cref{thm:b3:modules:frobenius}）；
也可比较秩：$\operatorname{rk} u = 2 \neq 1 =
\operatorname{rk} v$。对 $n \leq 3$：$\chi$ 与 $\mu$ 决定
对每个特征值 $\lambda$（在分裂域上），
对应块的总大小 $m_\lambda \leq 3$，$\lambda$ 块中最大者为
$r_\lambda$；而 $m \leq 3$ 的划分由其
最大部分决定（$m = 3, r = 2$ 必为 $(2,1)$，等等）。
故\omterm{thm:b3:modules:structure}{初等因子}相同，
并由\cref{cor:b3:modules:descent}，将相似性
下降到基域上。
\end{solution}

\begin{solution}{exo:b3:modules:9}
(i)$\Rightarrow$(ii)：若 $V = K[u]x$，则 $V \cong
K[X]/\operatorname{Ann}(x)$，且 $\operatorname{Ann}(x) = (\mu_u)$
（多项式零化 $x$，当且仅当它零化整个 $V = K[u]x$，因为
$P(u)Q(u)x = Q(u)P(u)x$）。故 $n = \dim V = \deg \mu_u$；
由 $\mu_u \mid \chi_u$ 且 $\deg\chi_u = n$，首一性给出 $\mu_u =
\chi_u$。

(ii)$\Rightarrow$(iii)：$\deg\chi_u = \sum_i \deg P_i$ 且
$\mu_u = P_s$（\cref{thm:b3:modules:frobenius}）；次数相等
迫使 $s = 1$。

(iii)$\Rightarrow$(i)：$V \cong K[X]/(P_1)$ 是循环的，
由 $\bar 1$ 的原像生成。

\omterm{def:b3:modules:companion}{友矩阵}即 $V = K[X]/(P)$ 本身的情形：$x = \bar
1$，即 $e_1$，是循环向量。对于对角矩阵
$\operatorname{diag}(\lambda_1, \dots, \lambda_n)$：$\chi = \prod
(X - \lambda_i)$，$\mu = \prod_{\lambda \text{ 互不相同}} (X -
\lambda)$；二者相等当且仅当 $\lambda_i$ 两两不同：
对角矩阵有循环向量，当且仅当其对角元
两两不同（此时 $x = (1, \dots, 1)$ 可取，
由范德蒙德行列式）。
\end{solution}

\begin{solution}{exo:b3:modules:10}
史密斯标准形：$M = Q\operatorname{diag}(d_1,\dots,d_n)P$；
由\cref{exo:b3:modules:4} 得 $[\Z^n : M\Z^n] = \prod\abs{d_i} =
\abs{\det M}$。若 $\det M = \pm 1$：伴随矩阵公式
$M^{-1} = (\det M)^{-1}\,{}^{t}\!\operatorname{com}(M)$ 的元素为
整数，故 $M \in GL_n(\Z)$；反之，$MM^{-1} = I$
给出 $\det M \cdot \det M^{-1} = 1$（在 $\Z$ 中），故 $\det M = \pm1$。

指数为 $m$ 的 $\Z^2$ 子群：此类子群 $L$ 的秩为
$2$（有限指数），且有唯一的\emph{埃尔米特
标准形}基 $\bigl(\begin{smallmatrix} a & b\\ 0 &
d\end{smallmatrix}\bigr)$：$d$ 由 $L \cap (\{0\}
\times \Z) = \{0\} \times d\Z$，$a$ 由 $\pi_1(L) = a\Z$（第一
坐标）刻画，而 $b$ 随后在模 $d$ 意义下唯一；规范化取 $a, d
> 0$ 且 $0 \leq b < d$。指数为 $ad = m$。计数：对每个
对每个因数 $d \mid m$（$a = m/d$），有 $d$ 个
可取的 $b$ 值；总数为 $\sum_{d \mid m} d = \sigma_1(m)$。
\end{solution}

\begin{solution}{exo:b3:modules:11}
(a) 由结构定理，$G \cong \prod_i\Z/d_i\Z$，且
$x^k = e$ 按坐标分解。在 $\Z/d\Z$ 中：$kx \equiv 0
\pmod d$ 恰有 $\gcd(k, d)$ 个解（$x$ 必为
$d/\gcd(k,d)$ 的倍数，这样的倍数共 $\gcd(k, d)$ 个）。
再对各因子相乘。

(b) 若 $G = \Z/d_s\Z$ 为循环群（$s = 1$），则计数为
$\gcd(k, d_s) \leq k$。若 $s \geq 2$：取 $k = d_1$；
计数为 $\prod_i\gcd(d_1, d_i) = d_1^{\,s} > d_1$（每个
$\gcd$ 都等于 $d_1$，由整除链）：方程
$x^{d_1} = e$ 的解多于 $d_1$ 个。

(c) 在域中，$X^k - 1$ 至多有 $k$ 个根
（\cref{ch:b3:rings}：整环上次数为 $k$ 的非零多项式），
故每个有限子群 $G \leq K^\times$
满足 (b) 的判别准则：$G$ 为循环群 —— 这给出
循环性的一行结构性证明，即
$\mathbb F_q^\times$ 的循环性，它
与\cref{ch:b3:galois} 中的计数证明互为补充。
\end{solution}

\begin{solution}{exo:b3:modules:12}
(a) 若 $MN = I$ 且 $N$ 为整数矩阵：$\det M\det N = 1$，
两者均为整数，故 $\det M = \pm1$。反之，若 $\det M =
\pm1$，则余因子公式 $M^{-1} =
\frac1{\det M}\,{}^t\!\operatorname{com}(M)$ 的元素为
整数。

(b) 左乘 $E^{-k}$ 将 $k$ 倍的
第 $2$ 行从第 $1$ 行中减去；左乘 $F^{-k}$ 则将 $k$ 倍的第 $1$ 行
从第 $2$ 行中减去。设 $M = \bigl(\begin{smallmatrix}a & *\\ c &
*\end{smallmatrix}\bigr) \in SL_2(\Z)$，则第一列
$(a, c)$ 是幺模向量（$\gcd(a, c) = 1$：它
整除 $\det M = 1$）。对 $(a, c)$ 用这些行变换施行欧几里得算法：
经有限步后，该列变为
$(\pm1, 0)$。此时矩阵为 $\pm
\bigl(\begin{smallmatrix}1 & b\\ 0 & 1\end{smallmatrix}
\bigr) = \pm E^{b}$（行列式保持为 $1$）。剩下要
将 $-I$ 用生成元表示：$(EF^{-1}E)^2 =
\bigl(\begin{smallmatrix}0 & 1\\ -1 &
0\end{smallmatrix}\bigr)^2 = -I$（验证旋转矩阵的
平方）。逆向回溯可知，$M$ 是 $E^{\pm1},
F^{\pm1}$ 的一个词。

(c) 史密斯算法（\cref{met:b3:modules:smith}）正是使用
这类行与列的变换（另加交换与变号，
二者在相差行列式符号的意义下本身也是初等运算之积）：
在 $\Z$ 上，每个矩阵都可写成
$U\,D\,V$，其中 $U, V$ 为初等矩阵之积，且
$D$ 为史密斯形——（b）是 $2\times2$、行列式为 $1$ 的
一般事实的实例：$E$ 型矩阵生成
$SL_n(\Z)$。
\end{solution}

\begin{solution}{pb:b3:modules:1}
\textbf{1.} $\varphi$ 可加，且 $K[X]$-线性：$\varphi(X
\cdot (Q_i)_i) = \sum_i (XQ_i)(M)e_i = M\sum_i Q_i(M)e_i = X
\cdot \varphi\bigl((Q_i)_i\bigr)$，$V_M$ 的模结构
为 $X \cdot v = Mv$。映射是满射：常值向量给出
整个 $K^n$。第 $j$ 列（取自 $XI - M$）为 $Xe_j - \sum_i
m_{ij}e_i$，其像为 $Me_j - \sum_i m_{ij}e_i = 0$。

\textbf{2.} 模去各列，$Xe_j \equiv \sum_i m_{ij}e_i$：
任一多项式向量，按最高次数作归纳，可约化为
\emph{常值}向量 $c \in K^n$。若原向量
属于 $\ker\varphi$，则 $\varphi(c) = c = 0$（在常值上，$\varphi$ 即是
$K^n = V_M$ 的同一视），故该
向量属于列张成的子模：$\ker\varphi = (XI_n -
M)K[X]^n$。因此 $V_M \cong K[X]^n/(XI - M)K[X]^n$，且
在 $K[X]$ 上的史密斯算法（\omterm{thm:b3:modules:smith}{不变因子}全非零，其积为
$\det(XI - M) = \chi_M$）给出 $V_M \cong \bigoplus_i
K[X]/(f_i)$：非常值的 $f_i$ 为
\omterm{thm:b3:modules:frobenius}{相似不变量}，可由子式的最大公因式之商算出，
（\cref{thm:b3:modules:smith}）。

\textbf{3.} $\lambda I_n$：$XI - \lambda I$ 本身已是史密斯形：
\omterm{thm:b3:modules:smith}{不变因子}为 $(X - \lambda, \dots, X - \lambda)$，共 $n$ 个。
对角元两两不同时：$V \cong \bigoplus_i K[X]/(X -
\lambda_i)$，各模两两互素（\omterm{def:b3:rings:ideal}{理想}），由中国剩余定理合并为
单一的循环模 $K[X]/\bigl(\prod_i(X - \lambda_i)\bigr)$：
其唯一\omterm{thm:b3:modules:smith}{不变因子}为 $\chi$。$J_n(0)$：$\mu = X^n = \chi$ 迫使
只有一个\omterm{thm:b3:modules:smith}{不变因子} $X^n$。$\operatorname{diag}(J_2(0), J_1(0))$：
\omterm{thm:b3:modules:structure}{初等因子}为 $X^2, X$：\omterm{thm:b3:modules:smith}{不变因子} $P_1 = X \mid P_2 =
X^2$。

\textbf{4.} $P \mapsto P(u)$ 将 $K[X]$ 满射至 $K[u]$，
其核为 $(\mu_u)$（由极小多项式的定义），故 $K[u]
\cong K[X]/(\mu_u)$，维数为 $\deg\mu_u = \deg P_s = n_s$。

\textbf{5.} 由 $K[X]$-线性可得 $f(\bar Q) = f(Q\cdot\bar 1)
= Q\,f(\bar 1)$。类 $\bar 1$ 满足 $P \bar 1 = 0$，故
$P f(\bar 1) = 0$ 是必要的。反之若 $P\bar c = 0$，
则 $f(\bar Q) = Q\bar c$ 良定义（$Q \equiv Q' \bmod P
\Rightarrow (Q - Q')\bar c \in$ 为 $P\bar c = 0$ 的倍式）且
$K[X]$-线性。

\textbf{6.} 设 $g = \gcd(P, Q)$，$P = gP'$，$Q = gQ'$，其中
$\gcd(P', Q') = 1$。在 $K[X]/(Q)$ 中：$P\bar c = 0 \iff Q \mid Pc
\iff Q' \mid P'c \iff Q' \mid c$（由欧几里得引理，$\gcd(P', Q') = 1$）。
故容许的 $\bar c$ 构成由
$\bar{Q'}$ 生成的子模，其零化子为 $\{R : Q \mid RQ'\} = (g)$：
该子模 $\cong K[X]/(g)$。结合第 5 题，
$\operatorname{Hom}(K[X]/(P), K[X]/(Q)) \cong K[X]/(\gcd(P,Q))$，
维数为 $\deg\gcd(P,Q)$。

\textbf{7.} $v$ 与 $u$ 可交换，当且仅当 $v$ 与每个
$P(u)$ 可交换，即 $v$ 为 $K[X]$-线性：$\mathcal C(u) =
\operatorname{End}_{K[X]}(V)$。把 $V =
\bigoplus_j V_j$ 上的态射写成矩阵 $(f_{ij})$，$f_{ij} \in
\operatorname{Hom}(V_j, V_i)$（与单射和
投影复合），由第 6 题得
\[
\dim \mathcal C(u) = \sum_{i,j} \deg\gcd(P_i, P_j)
= \sum_{i,j} n_{\min(i,j)}
= \sum_{k=1}^{s} \bigl(2(s - k) + 1\bigr)\,n_k ,
\]
利用整除链（$\gcd(P_i, P_j) =
P_{\min(i,j)}$），并在最后一步用到 $\min(i,j) = k$
恰在 $2(s - k) + 1$ 个数对 $(i, j)$ 上成立。

\textbf{8.} 由 $2(s-k)+1 \geq 1$，$\dim\mathcal C(u) \geq
\sum_k n_k = n$，等号成立当且仅当 $s = 1$，即当且仅当 $u$ 为
循环的（\cref{exo:b3:modules:9}）。对 $u = \lambda\,\mathrm{id}$：
$s = n$，所有 $n_k = 1$：$\dim = \sum_{k=1}^n (2(n-k)+1) = n^2$
——正确，因 $\mathcal C(\lambda\,\mathrm{id}) = \mathcal
L(V)$。

\textbf{9.} 由公式：\omterm{thm:b3:modules:smith}{不变因子}为 $(X, X^2)$，故 $s = 2$，$n_1 =
1$，$n_2 = 2$，$\dim = 3\cdot1 + 1\cdot2 = 5$。直接计算：在
基 $(e_1, e_2, e_3)$ 中，$u e_2 = e_1$，$ue_1 = ue_3 = 0$，
由 $Au = uA$（$A = (a_{ij})$）得条件
$a_{21} = a_{23} = a_{31} = 0$ 与 $a_{11} = a_{22}$，于是有五个自由
参数 $a_{11}{=}a_{22}, a_{12}, a_{13}, a_{32}, a_{33}$。

\textbf{10.} 多项式 $P(u)$ 与任一与 $u$ 可交换的
元素可交换（它是 $u$ 的幂之和）：$K[u]
\subseteq \mathcal C(\mathcal C(u))$。又 $u \in \mathcal C(u)$，
故任一 $w \in \mathcal C(\mathcal C(u))$ 与 $u$ 可交换。

\textbf{11.} 设 $V = K[u]x$，$v \in \mathcal C(u)$。记
$v(x) = P(u)x$（循环性）。对任意 $y = Q(u)x$：$v(y) =
vQ(u)x = Q(u)v(x) = Q(u)P(u)x = P(u)y$。故 $v = P(u)$，$\mathcal
C(u) = K[u]$。于是 $\mathcal C(\mathcal C(u)) = \mathcal
C(K[u]) = \mathcal C(u) = K[u]$（与 $K[u]$ 全体可交换
等价于与 $u$ 可交换）。定理在
循环情形成立。

\textbf{12.} $\pi_i$ 为 $K[X]$-线性（分解是
子模的直和），故 $\pi_i \in \mathcal C(u)$，且 $w$
与之可交换：$w(V_i) = w\pi_i(V) = \pi_i w(V) \subseteq
V_i$。限制 $w_i = w\restriction_{V_i}$ 与
循环的 $u_i = u\restriction_{V_i}$ 可交换（第 10 题），且 $w_i
\in \mathcal C(u_i) = K[u_i]$（第 11 题）：$w_i = Q_i(u_i) =
Q_i(u)\restriction_{V_i}$，其中 $Q_i \in K[X]$。

\textbf{13.} $\eta_{ij}(R(u)x_j) = R(u)x_i$ 的良定性：
若 $R(u)x_j = 0$，则 $P_j \mid R$，又 $P_i \mid P_j$，得
$P_i \mid R$，故 $R(u)x_i = 0$（$\operatorname{Ann}(x_i) =
(P_i)$）。由构造，$\eta_{ij}$ 是 $K[X]$-线性的，且
$\tilde\eta_{ij} = \eta_{ij}\pi_j$ 也是
$K[X]$-态射 $V \to V$ 的复合：$\tilde\eta_{ij} \in \mathcal
C(u)$。

\textbf{14.} 将 $w\tilde\eta_{ij} = \tilde\eta_{ij}w$ 在
$x_j$ 上取值：左边为 $w(x_i) = Q_i(u)x_i$；右边为
$\eta_{ij}\bigl(Q_j(u)x_j\bigr) = Q_j(u)x_i$。于是 $(Q_i -
Q_j)(u)\,x_i = 0$：$P_i \mid Q_i - Q_j$ 对一切 $i \leq j$ 成立。特别地，
取 $Q = Q_s$，则 $Q \equiv Q_i \pmod{P_i}$，
$Q(u)$ 与 $Q_i(u)$ 在 $V_i$ 上一致（$P_i(u)$ 将其零化）。因此 $w
= Q(u)$ 在每个 $V_i$ 上成立，从而在 $V$ 上亦然：$\mathcal C(\mathcal C(u))
\subseteq K[u]$，结合第 10 题，$\mathcal C(\mathcal
C(u)) = K[u]$。

\textbf{15.} 第一个断言即第 10--14 题（或对循环的
$u$，仅用第 11 题）。对 $u = \mathrm{id}$，$n \geq
2$：$\mathcal C(u) = \mathcal L(V)$ 的维数为 $n^2$，而
$\dim K[u] = \deg\mu_u = 1$。因此 $\mathcal C(u) = K[u]$ 严重不成立；
然而双\omterm{ex:b3:groups:actions}{中心化子}定理
仍成立（$\mathcal C(\mathcal L(V)) = K\,\mathrm{id} = K[u]$：
矩阵代数的中心恰为数量矩阵）。循环性正是
使\emph{单个}\omterm{ex:b3:groups:actions}{中心化子}已是多项式；
而\emph{双重}\omterm{ex:b3:groups:actions}{中心化子}总是多项式。

\textbf{16.} 若 $P(XI - M)Q = S$ 为史密斯约化（$P, Q$
在 $K[X]$ 上可逆），转置得 ${}^tQ\,(XI -
{}^tM)\,{}^tP = {}^tS = S$：史密斯标准形相同，故 $XI - M$ 与
$XI - {}^tM$ 有相同的\omterm{thm:b3:modules:smith}{不变因子}，即 $M$ 与
${}^tM$ 有相同的\omterm{thm:b3:modules:frobenius}{相似不变量}
（见\cref{cor:b3:modules:descent}），故二者相似。

\textbf{17.} $M$ 在 $L$ 上的\omterm{thm:b3:modules:frobenius}{相似不变量}即
$XI - M$ 在 $L[X]$ 中的\omterm{thm:b3:modules:smith}{不变因子}。考虑
$XI - M$ 在 $K[X]$ 上的一个史密斯约化——即 $P, Q$ 为 $K[X]$ 上的可逆矩阵，
对角且满足整除链——\emph{同样}是
$L[X]$ 上有效的史密斯约化（$P, Q$ 仍可逆：其
行列式为非零常数），且首一\omterm{thm:b3:modules:smith}{不变因子}
唯一：在 $K$ 上算出的\omterm{thm:b3:modules:smith}{不变因子}与
在 $L$ 上算出的\omterm{thm:b3:modules:smith}{不变因子}相同。故 $M \sim_L N$ 当且仅当二者有相同的
\omterm{thm:b3:modules:smith}{不变因子}，亦即 $M \sim_K N$。特别地，
在 $\C$ 上相似的实矩阵在 $\R$ 上也相似——此结论
常以分析方法证明（取可逆矩阵 $P + \iu
Q$ 作特化），这里则从结构上得出。

\textbf{18.} 将 $u = \bigoplus J_{m_t}(0)$ 分解为
幂零若尔当块之直和。在大小为 $m$ 的一个块中，
$\operatorname{rk}J_m^k = \max(m - k, 0)$，故
$\operatorname{rk}J_m^{k-1} - \operatorname{rk}J_m^k = 1$，若 $m
\geq k$，否则为 $0$。对各块求和：
$\operatorname{rk}u^{k-1} - \operatorname{rk}u^k =
\#\{t : m_t \geq k\}$。因此秩序列确定
多重集 $(m_t)$——即 $n$ 的一个分拆——反之
每个分拆都可实现，幂零类 $\leftrightarrow$
$n$ 的分拆，对每个域都成立。$M_5$：$p(5) = 7$ 个类
（$5$；$4{+}1$；$3{+}2$；$3{+}1{+}1$；$2{+}2{+}1$；
$2{+}1{+}1{+}1$；$1^5$）。

\textbf{19.} (a) 在 $\Q$ 上（或任一特征为 $0$ 的域），
$D^n = 0$，$D^{n-1}(X^{n-1}) = (n-1)!\, \neq 0$：$D$ 是
在 $n$ 维空间上指数为 $n$ 的幂零算子，
从而是循环的，唯一\omterm{thm:b3:modules:smith}{不变因子}为 $X^n$（取 $x = X^{n-1}$，它生成：
其各阶导数张成全空间）。(b) 在 $\mathbb F_p$ 上，$p
< n$ 时：$D^p = 0$，因 $p$ 阶导数作用于每个
单项式 $X^m$ 时都含因子 $m(m-1)\cdots(m-p+1)$，
即 $p$ 个连续整数之积，故 $\equiv 0 \bmod p$。
记 $n = ap + r$，$0 \leq r < p$。于是 $\ker D^k$ 由
单项式 $X^m$ 张成，其中 $D^kX^m = 0$；
对 $m < n$ 的指数按模 $p$ 的剩余类计数：$\dim\ker D^k = ak + \min(r,
k)$，$0 \leq k \leq p$，故 $\operatorname{rk}D^{k-1} -
\operatorname{rk}D^k = \dim\ker D^k - \dim\ker D^{k-1} = a +
\mathbf 1_{k \leq r}$。由第 18 题，此分拆有 $a$
个大小恰为 $p$ 的块，以及（若 $r > 0$）一个大小为
$r$ 的块：\omterm{thm:b3:modules:frobenius}{相似不变量}为 $X^r \mid X^p \mid \dots \mid X^p$。
特征的不同改变了数学中最熟悉的
算子的标准形。

\textbf{20.} $M \in GL_2$ 有 $s \in \{1, 2\}$ 个
\omterm{thm:b3:modules:smith}{不变因子}。若 $s = 2$：$P_1 = P_2 = X - a$（$a \neq 0$，
即可逆性），亦即 $M = aI$，是中心元。若 $s = 1$：$M$ 是
循环的，其特征多项式 $=$ 极小多项式 $\chi$，
次数为 $2$，各类与可能的 $\chi$ 一一对应，
且 $\chi(0) \neq 0$：$\chi$ 分裂且有相异根 $a \neq
b$（\omterm{def:b3:modules:companion}{友矩阵} $\sim$ 对角矩阵）；$\chi = (X - a)^2$（\omterm{def:b3:modules:companion}{友矩阵}，
非半单）；$\chi$ 不可约。每种情形恰对应一个类——
\omterm{thm:b3:modules:smith}{不变因子}是完全不变量。

\textbf{21.} 中心情形：$q - 1$ 种 $a$ 的取法。两两相异的分裂
特征值：无序对 $\{a, b\} \subseteq
\mathbb F_q^\times$，$a \neq b$，共 $\binom{q-1}2$ 个类。
极小多项式为 $(X-a)^2$：$q - 1$ 个类。常数项非零的不可约二次式：
凡不可约二次式皆合格，
（其根均非零），共有 $\frac{q^2 - q}2$
个首一不可约二次式（即 $q^2$ 个首一二次式减去
$\binom q2 + q = \frac{q^2+q}2$ 个可分裂者）。总和为：
\[
(q - 1) + \frac{(q-1)(q-2)}2 + (q - 1) + \frac{q^2 - q}2
= q^2 - 1 .
\]

\textbf{22.} 若 $M$ 不循环，则 $s = 2$，且 $M$ 为数量矩阵
（第 20 题的二分法在 $M_2$ 中成立，无论是否可逆：两个
次数为 $1$ 的\omterm{thm:b3:modules:smith}{不变因子}，满足 $P_1 \mid P_2$ 且 $P_1 =
P_2$ 时必有 $M = aI$）。数量矩阵共 $q$ 个，在全部 $q^4$
个矩阵中：循环矩阵的概率为 $1 - q^{-3}$。一般地，
$M_n$ 中非循环矩阵的轨迹，正是 $(n-1)\times(n-1)$ 阶子式
（取自 $XI - M$）有公因子之处——一个真正的代数条件
——故其占比为 $O(1/q)$ 阶的小量，对大的 $q$ 而言：
$\chi = \mu$ 的矩阵才是常态，而第三部分的粘合论证
正是为例外情形所付出的代价。

\textbf{23.} $\mathcal C(u)$ 中心的元素
属于 $\mathcal C(u)$，并与其中每个元素交换，
$\mathcal C(u)$ 中的元素即属于 $\mathcal C(\mathcal C(u)) =
K[u]$（第 14 题）。反之 $K[u] \subseteq \mathcal C(u)$，
且每个 $P(u)$ 都与每个 $v \in \mathcal C(u)$ 交换，
（这样的 $v$ 与 $u$ 交换，从而与 $u$ 的每个幂交换）：
$K[u]$ 位于 $\mathcal C(u)$ 的中心，因此 $Z(\mathcal C(u))
= K[u]$。由此，$\mathcal C(u)$ 可交换当且仅当
$\mathcal C(u) = Z(\mathcal C(u)) = K[u]$；此时
$\dim\mathcal C(u) = \dim K[u] = n_s \leq n$，而第 8 题
给出 $\dim\mathcal C(u) \geq n$，于是 $n_s = n$ 且 $s = 1$，
即 $u$ 为循环的。反之，对循环的 $u$，由第 11 题
得 $\mathcal C(u) = K[u]$，可交换。从结构上看，
与自身中心相等的矩阵代数恰为
多项式代数 $K[u]$，其中 $u$ 循环。

\textbf{24.} 弗罗贝尼乌斯公式中的每个系数 $2s - 2i + 1$
均为奇数，故
\[
\dim\mathcal C(u) = \sum_{i=1}^s(2s - 2i + 1)\,n_i
\equiv \sum_{i=1}^s n_i = n \pmod 2 .
\]
当 $n = 4$ 时，\omterm{thm:b3:modules:smith}{不变因子}的可能次数序列 $n_1 \leq \dots
\leq n_s$，其和为 $4$，共有
$(4)$、$(1, 3)$、$(2, 2)$、$(1, 1, 2)$、$(1, 1, 1, 1)$；每一种都可实现，
例如取 $P_i = X^{n_i}$ 的幂零矩阵即可实现，
（整除链自动成立）。由公式算得
依次为
\[
1\cdot4 = 4, \quad 3 + 3 = 6, \quad 6 + 2 = 8, \quad
5 + 3 + 2 = 10, \quad 7 + 5 + 3 + 1 = 16 .
\]
故取值集为 $\{4, 6, 8, 10, 16\}$，全为
奇偶性正确的偶数，但 $12$ 与 $14$ 从不出现——在
近循环序列与标量情形的 $n^2$ 之间存在一段空隙。

\textbf{25.} $\abs{GL_2(\mathbb F_3)} = (q^2 - 1)(q^2 - q) =
8 \cdot 6 = 48$。中心型：$I$ 与 $2I$，两个大小为
$1$ 的类。其余三种类型中 $M$ 循环（由第
20 题），故其在 $GL_2$ 中的\omterm{ex:b3:groups:actions}{中心化子}为
$\mathcal C(M) = K[M]$ 中的可逆元所成之群（第 11 题），类的大小
$=48/\abs{K[M]^\times}$（由\omterm{def:b3:groups:action}{轨道}–\omterm{def:b3:groups:action}{稳定子}定理）。不同的分裂特征值
仅 $\{1, 2\}$ 这一对，共一个类；$K[M]
\cong \mathbb F_3 \times \mathbb F_3$（由中国剩余定理，对 $\chi = (X-1)
(X-2)$），可逆元个数 $2 \cdot 2 = 4$，大小为 $48/4 = 12$。极小
多项式 $(X - a)^2$，$a \in \{1, 2\}$：两个类；$K[M] \cong
\mathbb F_3[X]/((X-a)^2)$，可逆元个数 $q^2 - q = 6$，
可逆元的常数项 $\neq 0$（中心化之后），大小为 $48/6 = 8$。
不可约的 $\chi$：首一不可约二次多项式在
$\mathbb F_3$ 上共 $(9 - 3)/2 = 3$ 个，即
\[
X^2 + 1, \qquad X^2 + X + 2, \qquad X^2 + 2X + 2,
\]
（在 $\mathbb F_3$ 中无根：逐一验算 $0, 1, 2$）；共三个类，
$K[M] \cong \mathbb F_9$，可逆元个数 $q^2 - 1 = 8$，大小为 $48/8 =
6$。\omterm{cor:b3:groups:classeq}{类方程}：$2\cdot1 + 1\cdot12 + 2\cdot8 + 3\cdot6 =
2 + 12 + 16 + 18 = 48$；又 $2 + 1 + 2 + 3 = 8 = q^2 - 1$
个\omterm{ex:b3:groups:actions}{共轭类}，与第 21 题相符。
\end{solution}
